\exercisesection{3}

\begin{enumerate}
\def\labelenumi{\arabic{enumi}.}
\tightlist
\item
  \begin{enumerate}
  \def\labelenumii{(\alph{enumii})}
  \tightlist
  \item
    Let A be a ring which is not necessarily commutative and \(I_{s}\) (resp. \(I_{d}\) ) the set of elements of A which have no left (resp. right) inverse. Show that the following conditions are equivalent: (1) the sum of two elements of \(I_{s}\) belongs to \(I_{s}\) ; (2) \(I_{s}\) is a left ideal; (3) there is a greatest left ideal J in the set of left ideals \(\neq A\) . If these conditions hold, the opposite ring \(A^{0}\) also satisfies these conditions, \(I_{s} = I_{d} = J\) , J is the unique maximal (left or right) ideal of A and hence the Jacobson radical of A, every element \(x \notin J\) is invertible and A/J is a field (cf Algebra, Chapter VIII, §6, no. 3). Then we also say that A is a local ring.
  \end{enumerate}
\end{enumerate}

\begin{enumerate}
\def\labelenumi{(\alph{enumi})}
\setcounter{enumi}{1}
\tightlist
\item
  Show that in a local ring there exists no idempotent other than 0 and 1.
\end{enumerate}

\begin{enumerate}
\def\labelenumi{\arabic{enumi}.}
\setcounter{enumi}{1}
\tightlist
\item
  \begin{enumerate}
  \def\labelenumii{(\alph{enumii})}
  \tightlist
  \item
    Let A be a (commutative) local ring whose maximal ideal m is principal and such that \(\bigcap_{n\geqslant1}m^{n}=0\) (cf.~Chapter III, §3, no. 2, Corollary to Proposition 4). Show that the only ideals \(\neq0\) and \(\neq A\) of A are the powers \(m^{n}\) ; deduce that A is Noetherian and is *either a discrete valuation ring (Chapter VI)* or a quasi-principal ring in which 1 is an indecomposable idempotent (Algebra, Chapter VII, §1, Exercise 6).
  \end{enumerate}
\end{enumerate}

\begin{enumerate}
\def\labelenumi{(\alph{enumi})}
\setcounter{enumi}{1}
\tightlist
\item
  Let A be the ring of germs at the point t = 0 of real-valued functions defined and continuous in a neighbourhood of 0 and differentiable at the point 0 (General Topology, Chapter I, 3rd ed., § 6, no. 10). Show that A is a local ring whose maximal ideal m is generated by the germ of the function \(j: t \mapsto t; m^{n} \neq 0\) for all n but A is not a principal ideal domain. If c is the germ of the function \(t \mapsto \exp(-1/t^{2})\) and p a prime ideal of A not containing c, show that the quotient ring B = A/p is a local integral domain which is not a principal ideal domain and in which the maximal ideal is principal.
\end{enumerate}

¶ 3. Let A be a local ring (not necessarily commutative, cf.~Exercise 1) and m its Jacobson radical.

\begin{enumerate}
\def\labelenumi{(\alph{enumi})}
\item
  Let L be a free left A-module and \(x \neq 0\) an element of L. Let \((e_{\iota})_{\iota \in I}\) be a basis of L for which the number of components \(\neq 0\) of x is the least possible; if \(x = \sum_{\iota \in I} \xi_{\iota} e_{\iota}\) , let J be the set of \(\iota \in I\) such that \(\xi_{\iota} \neq 0\) . Show that none of the \(\xi_{\iota} (\iota \in J)\) belongs to the right ideal of A generated by the others.
\item
  With the hypotheses and notation of (a), let P, Q be two complementary submodules of L such that \(x \in P\) ; for all \(\iota \in I\) , let \(e_{\iota} = y_{\iota} + z_{\iota}\) , where \(y_{\iota} \in P\) , \(z_{\iota} \in Q\) . Show that the family of \(y_{\iota}\) for \(\iota \in J\) and \(e_{\iota}\) for \(\iota \notin J\) is a basis of L. (Prove, using (a), that, if \(z_{\kappa} = \sum_{\iota \in I} \zeta_{\kappa_{\iota}} e_{\iota}\) , the components \(\zeta_{\kappa_{\iota}}\) such that \(\kappa \in J\) and \(\iota \in J\) necessarily belong to m and use Corollary 2 to Proposition 4.) Deduce that there exists a free submodule of P which contains x and is a direct factor of P.
\item
  Show that every projective A-module P is free (using Kaplansky's Theorem (Algebra, Chapter II, § 2, Exercise 4), reduce it to the case where P is generated by a countable family of elements and apply (b) to this case).
\item
  Give an example of a local ring A and a flat A-module M which is not free and satisfies mM = M (take \(A = \mathbf{Z}_{(p)}\) ). *(In Chapter III, § 3, we shall give examples of faithfully flat A-modules which are not free over a Noetherian local ring A.)*
\item
  Let A be a local ring and M a finitely generated flat A-module. Show that M is a free A-module (use Exercise 23 (e) of Chapter I, § 2).
\end{enumerate}

\begin{enumerate}
\def\labelenumi{\arabic{enumi}.}
\setcounter{enumi}{3}
\item
  Let A be a local ring (not necessarily commutative) and m its Jacobson radical. Let M and N be two free A-modules (finitely generated or not) and \(u: M \rightarrow N\) a homomorphism such that \(1 \otimes u: M/mM \rightarrow N/mN\) is bijective. Show that u is bijective. (Prove first that u is surjective and then that u is injective, reducing it to the case where M and N are finitely generated and using respectively Corollary 1 to Proposition 4 and Proposition 6.)
\item
  Show by an example that Proposition 7 does not extend to the case where the local ring A is not reduced.
\item
  \begin{enumerate}
  \def\labelenumii{(\alph{enumii})}
  \tightlist
  \item
    Let M be the Z-module defined in Algebra, Chapter VII, §3, Exercise 22 (h); let T be the torsion submodule of M. Show that, for every prime number p, the submodule \(\mathrm{T}_{(p)}\) of \(\mathrm{M}_{(p)}\) is a direct factor of \(\mathrm{M}_{(p)}\) , although T is not a direct factor of M.
  \end{enumerate}
\end{enumerate}

\begin{enumerate}
\def\labelenumi{(\alph{enumi})}
\setcounter{enumi}{1}
\tightlist
\item
  Let \(\mathbf{M}\) be the \(\mathbf{Z}\) -module defined in Algebra, Chapter VII, § 3, Exercise 5; show that, for every prime number \(p\) , \(\mathrm{M}_{(p)}\) is a free \(\mathbf{Z}_{(p)}\) -module, although \(\mathbf{M}\) is not a free \(\mathbf{Z}\) -module.
\end{enumerate}

\begin{enumerate}
\def\labelenumi{\arabic{enumi}.}
\setcounter{enumi}{6}
\item
  Let A be a ring and \((\mathrm{S}_{\lambda})_{\lambda \in L}\) a family of multiplicative subsets of A such that, for every maximal ideal m of A, there exists \(\lambda \in L\) such that \(m \cap S_{\lambda} = \varnothing\) (cf.~§2, Exercise 8). Let M, N be two A-modules and u: M → N a homomorphism. For u to be surjective (resp. injective, bijective, zero), it is necessary and sufficient that, for all \(\lambda \in L\) , \(S_{\lambda}^{-1}u\) : \(S_{\lambda}^{-1}M \to S_{\lambda}^{-1}N\) be so. How can the Corollaries to Theorem 1 of no. 3 be generalized?
\item
  Let A be a ring, B an A-algebra (not necessarily commutative) and E a finitely generated left B-module.
\end{enumerate}

\begin{enumerate}
\def\labelenumi{(\alph{enumi})}
\item
  Let S be a multiplicative subset of A and let \(S^{-1}B = S^{-1}A \otimes_{A} B\) be given its \(S^{-1}A\) -algebra structure; \(S^{-1}E\) can then be considered as a left \(S^{-1}B\) -module (Algebra, Chapter VIII, § 7, no. 1) isomorphic to \(S^{-1}B \otimes_{B} E\) . Show that, if E is a flat B-module, \(S^{-1}E\) is a flat \(S^{-1}B\) -module.
\item
  Let \((\mathbf{S}_{\lambda})_{\lambda \in L}\) be a family of multiplicative subsets of A such that, for every maximal ideal m of A, there exists \(\lambda \in L\) such that \(m \cap S_{\lambda} = \varnothing\) . Show that \(\bigoplus_{\lambda \in L} S_{\lambda}^{-1} B\) is a faithfully flat B-module (left or right; cf.~§ 2, Exercise 8). Show that, if, for all \(\lambda \in L\) , \(S_{\lambda}^{-1} E\) is a flat (resp. faithfully flat) \(S_{\lambda}^{-1} B\) -module, then E is a flat (resp. faithfully flat) B-module (use Chapter I, § 3, Exercise 9, with \(C_{\lambda} = S_{\lambda}^{-1} A\) ).
\item
  Suppose that the \(S_{\lambda}\) satisfy condition (b) and that L is finite. Show that, if, for every \(\lambda \in L\) , \(S_{\lambda}^{-1}E\) is a finitely generated projective \(S_{\lambda}^{-1}B\) -module, then E is a finitely generated projective B-module.
\end{enumerate}

\begin{enumerate}
\def\labelenumi{\arabic{enumi}.}
\setcounter{enumi}{8}
\item
  Show that, for a (commutative) ring A to be absolutely flat (Chapter I, § 2, Exercise 17), it is necessary and sufficient that, for every maximal ideal m of A, \(A_{m}\) be a field (note that an absolutely flat local ring is necessarily a field and use the Corollary to Proposition 15).
\item
  Let A, B be two rings, \(\rho: A \rightarrow B\) a homomorphism, q a prime ideal of B and \(p = \bar{\rho}^{-1}(q)\) . Suppose that there exists a B-module N such that \(N_{q}\) is a flat A-module which is not reduced to 0 and \(N_{q}\) is a finitely generated \(A_{p}\) -module; then the prime ideal p is minimal in A. (Observe that \(N_{p}\) is then a faithfully flat \(A_{p}\) -module and use Corollary 4 to Proposition 11 of §2, no. 5.)
\item
  Let A be a ring, a an ideal of A and S the multiplicative subset of A consisting of the elements whose canonical images in A/a are not divisors of 0; let \(\Phi\) be the set of maximal elements of the set of ideals of A not meeting S (such that the ideals \(p \in \Phi\) are prime, cf.~§2, no. 5, Proposition 11). Show that a is the intersection of its saturations with respect to the ideals \(p \in \Phi\) (reduce it to the case where a = 0 and use Corollary 1 to Theorem 1).
\end{enumerate}

¶ 12. Let \(A = \prod_{i=1}^{n} A_i\) be the product of a finite family of local rings \(A_i\) , canonically identified with ideals of \(A\) . Let \(B\) be a subring of \(A\) such that \(\text{pr}_i B = A_i\) for \(1 \leqslant i \leqslant n\) . Show that \(B\) is the direct composition of at most \(n\) local rings and can only be the direct composition of \(n\) local rings if \(B = A\) . (Proceed by induction on \(n\) , considering the ideals \(a_i = B \cap A_i\) of \(B\) . Examine successively the case where \(a_i = A_i\) for at least one index \(i\) and the case where \(a_i \neq A_i\) for all \(i\) . Note that, if \(a_i \neq A_i\) , every maximal ideal of \(B\) contains \(a_i\) , and considering the ring \(B / a_i\) , conclude that there are at most \(n - 1\) distinct maximal ideals in \(B\) ; if then \(B\) is not a local ring, reduce it to the case where \(a_j = A_j\) for some \(j \neq i\) using the induction hypothesis and Exercise 1 (b).) Give an example where \(n > 1\) and \(B\) is a local ring but not isomorphic to any \(A_t\) (take the \(A_t\) to be algebras over the same field \(K\) , whose maximal ideals \(m_t\) are of zero square).

¶ 13. Let G be a finite group whose order n is a power \(p^{f}\) of a prime number p and K a field of characteristic p.

\begin{enumerate}
\def\labelenumi{(\alph{enumi})}
\item
  Let E be a finite set on which G operates (Algebra, Chapter I, § 7, no. 2); if \(E'\) is the set of elements of E which are invariant for all \(g \in G\) , show that \(\text{Card}(E) \equiv \text{Card}(E') \pmod{p}\) . Deduce that, if G is not reduced to its identity element e, the centre Z of G is not reduced to e (make G operate on itself by inner automorphisms). Conclude that G is solvable (Algebra, Chapter I, § 6, Exercise 14).
\item
  Let V be a vector space over K and \(\rho\) a homomorphism of G to \(\mathbf{GL}(V)\) . Let \(V'\) be the set of \(x \in V\) which are invariant under the linear mappings \(\rho(g)\) , where \(g \in G\) . Show that, if \(V'\) is reduced to 0, V is necessarily reduced to 0. (If \(x \in V\) and \(x \neq 0\) , apply (a) to the additive subgroup E of V generated by the elements \(\rho(g).x\) , where \(g \in G\) .)
\item
  Let \(\mathbf{A} = \mathbf{K}^{(\mathbf{G})}\) be the algebra of the group \(\mathbf{G}\) over \(\mathbf{K}\) (Algebra, Chapter III, § 1) and let \(\mathbf{I}\) be the vector subspace of \(\mathbf{A}\) (over \(\mathbf{K}\) ) generated by the elements \(1 - g\) where \(g \in \mathbf{G}\) . Show that \(\mathbf{I}\) is the Jacobson radical of \(\mathbf{A}\) and that \(\mathbf{A} / \mathbf{I}\) is isomorphic to \(\mathbf{K}\) (use (b) and the definition of Jacobson radical); deduce that \(\mathbf{I}\) is nilpotent.
\item
  Let V be a finitely generated A-module and V' the set of elements x of V such that g.x = x for all \(g \in G\) . Show that the following inequalities hold
\end{enumerate}

\[
(*) \dim_ {\mathbb {K}} (\mathrm{V}) \leqslant n. \dim_ {\mathbb {K}} (\mathrm{V} / \mathrm{IV}); \quad (* *) \dim_ {\mathbb {K}} (\mathrm{V}) \leqslant n. \dim_ {\mathbb {K}} (\mathrm{V} ^ {\prime}).
\]

(To establish (*), use Corollary 2 to Proposition 4; deduce (**) by considering the dual of the vector space V over K.) For the two sides of the inequality (*) (resp. (**)) to be equal, it is necessary and sufficient that V be a free A-module (cf.~Corollary 1 to Proposition 5).

\begin{enumerate}
\def\labelenumi{\arabic{enumi}.}
\setcounter{enumi}{13}
\tightlist
\item
  Let A be a ring, which is not necessarily commutative and whose Jacobson radical m is such that A/m is a field, and M a finitely generated free A-module.
\end{enumerate}

\begin{enumerate}
\def\labelenumi{(\alph{enumi})}
\item
  Let \(\mathbf{M}'\) be another finitely generated free A-module and \(u: \mathbf{M} \to \mathbf{M}'\) a homomorphism such that \(u(\mathbf{M})\) is a direct factor of \(\mathbf{M}'\) ; the number \(\dim(\mathbf{M})\) is called the rank of \(u\) and denoted by \(\operatorname{rg}(u)\) . For such a homomorphism \(\operatorname{Ker}(u)\) is a direct factor of \(\mathbf{M}\) and \(\operatorname{rg}(u) = \dim(\mathbf{M}) - \dim(\operatorname{Ker}(u))\) ; for \(u\) to be injective (resp. surjective), it is necessary and sufficient that \(\operatorname{rg}(u) = \dim(\mathbf{M})\) (resp. \(\operatorname{rg}(u) = \dim(\mathbf{M}')\) ). The transpose \(^t u\) is such that \(^t u(\mathbf{M}'*)\) is a direct factor of \(\mathbf{M}^*\) and \(\operatorname{rg}(^t u) = \operatorname{rg}(u)\) .
\item
  Let \(n = \dim(\mathbf{M})\) . The direct factors of \(\mathbf{M}\) of dimension 1 (resp. \(n - 1\) ) are called lines (resp. hyperplanes). An automorphism \(u \in \mathbf{GL}(\mathbf{M})\) distinct from the identity is called a transvection if there exists a hyperplane \(\mathbf{H}\) of \(\mathbf{M}\) all of whose elements are invariant under \(u\) ; then \(u(x) = x + a\phi(x)\) , where \(\phi\) is a linear form on \(\mathbf{M}\) such that \(\mathbf{H} = \operatorname{Ker}(\phi)\) and \(a \in \mathbf{H}\) ; obtain the converse. If \(\mathbf{A}\) is commutative, show that every automorphism \(u \in \mathbf{GL}(\mathbf{M})\) of determinant 1 is a product of transvections (observe that, in the matrix of \(u\) with respect to any basis of \(\mathbf{M}\) , every column contains at least one invertible element of \(\mathbf{A}\) and a matrix of the form \(I + E_{ij} (i \neq j)\) is the matrix of a transvection).
\item
  Give an example of an automorphism \(u \in \mathbf{GL}(\mathbf{M})\) such that the kernel of \(1 - u\) is not a direct factor of \(\mathbf{M}\) (take \(\mathbf{A}\) to be the local ring \(\mathbf{K}[[\mathbf{X}]]\) of formal power series in one indeterminate over a field \(\mathbf{K}\) and \(\mathbf{M} = \mathbf{A}\) ).
\item
  Give an example of direct factors N, P of M (necessarily free) such that \(N + P\) and \(N \cap P\) are not direct factors of M. (Take A = K{[}X{]}/(X²), where K is a field and M = A².)
\end{enumerate}

¶ 15. Let A be a (commutative) ring and M an A-module.

\begin{enumerate}
\def\labelenumi{(\alph{enumi})}
\item
  For a submodule \(M'\) of M to be pure (Chapter I, §2, Exercise 24), it is necessary and sufficient that, for every maximal ideal m of A, \(M_{m}'\) be a pure sub- \(A_{m}\) -module of \(M_{m}\) (use Theorem 1 of no. 3).
\item
  Suppose that A is a local ring with maximal ideal m and M a finitely generated free A-module. For a finitely generated submodule \(M'\) of M to be pure, it is necessary and sufficient that it be a direct factor of M. (Using Corollary 1 to Proposition 5 of no. 2, reduce it to proving that, if \(M' \subset mM\) , \(M'\) can only be a pure submodule of M if \(M' = 0\) .)
\end{enumerate}

\begin{enumerate}
\def\labelenumi{\arabic{enumi}.}
\setcounter{enumi}{15}
\tightlist
\item
  Let \((\mathbf{A}_{\lambda}, f_{\mu \lambda})\) be a direct system of local rings, such that the \(f_{\mu \lambda}\) are local homomorphisms; let \(m_{\lambda}\) be the maximal ideal of \(\mathbf{A}_{\lambda}\) and \(\mathbf{K}_{\lambda} = \mathbf{A}_{\lambda} / m_{\lambda}\) . Then \(\mathbf{A} = \varinjlim \mathbf{A}_{\lambda}\) is a local ring whose maximal ideal is \(m = \varinjlim m_{\lambda}\) and residue field is \(\mathbf{K} = \varinjlim \mathbf{K}_{\lambda}\) . Moreover, if \(m_{\mu} = \mathbf{A}_{\mu} m_{\lambda}\) for \(\lambda < \mu\) , then \(m = Am_{\lambda}\) for all \(\lambda\) .
\end{enumerate}

